\begin{sgastatement}{Definition 5.6}
Let \(S\) be a scheme, \(G\) an \(S\)-reductive group, and \(P\) a
parabolic subgroup of \(G\).  The subgroup \(P\) is said to be
\emph{minimal} if every parabolic subgroup \(Q\) of \(G\) contained
in \(P\) is equal to \(P\).
\end{sgastatement}
\SGARefTarget{sga3:xxvi:def:5:6}

Notice that this notion is not, in general, stable under passage to
the fibers.

\begin{sgastatement}{\SGAInlineRefTarget{sga3:xxvi:corollary:5-7}Corollary 5.7}
Let \(S\) be a semilocal scheme and \(G\) an \(S\)-reductive group.
\begin{enumerate}[label=\textup{(\roman*)}]
\item Let
\[
 t,t'\in\Of\bigl(\operatorname{Dyn}(G)\bigr)(S).
\]
If \(G\) has a parabolic subgroup of type \(t\) and one of type
\(t'\), then it has a parabolic subgroup of type \(t\cap t'\).  In
particular, the set of \(t(P)\), as \(P\) ranges over the parabolic
subgroups of \(G\), has a least element \(t_{\min}\).
\item Every parabolic subgroup of \(G\) contains a minimal parabolic
subgroup.  A parabolic subgroup of \(G\) is minimal if and only if it
has type \(t_{\min}\).  Any two minimal parabolic subgroups of \(G\)
are conjugate by an element of \(G(S)\).
\end{enumerate}
\end{sgastatement}
\SGARefTarget{sga3:xxvi:cor:5:7}

This follows immediately from \SGARefLink{sga3:xxvi:cor:5:4}{5.4} and \SGARefLink{sga3:xxvi:corollary:5-5}{5.5}\textup{(i)}.

\begin{sgastatement}{Remark 5.8}
A parabolic subgroup opposite to a minimal parabolic subgroup is
also minimal.  Consequently,
\[
 s(t_{\min})=t_{\min}.
\]
\end{sgastatement}
\SGARefTarget{sga3:xxvi:rem:5:8}

\begin{sgastatement}{\SGAInlineRefTarget{sga3:xxvi:corollary:5-9}Corollary 5.9}
Let \(S\) be a scheme, \(G\) an \(S\)-reductive group, and \(P\) a
parabolic subgroup of \(G\).  The canonical morphism
\[
 G\longrightarrow G/P=X
\]
makes \(G\) a Zariski-locally trivial \(P_X\)-bundle over \(X\).
If \(L\) is a Levi subgroup of \(P\), the canonical morphism
\textup{(cf.\ \SGARefLink{sga3:xxvi:cor:3:12}{3.12})}
\[
 G\longrightarrow G/L=Y
\]
makes \(G\) a Zariski-locally trivial \(L_Y\)-bundle over \(Y\).
\end{sgastatement}
\SGARefTarget{sga3:xxvi:cor:5:9}

It is enough to prove the following lifting assertion.  Given a
morphism \(S'\to S\), with \(S'\) local, every morphism \(S'\to X\)
\textup{(respectively \(S'\to Y\))} lifts to a morphism \(S'\to G\).
In other words, we may assume that \(S\) is local, and must show that
\[
 G(S)\longrightarrow X(S),
 \qquad\text{respectively}\qquad
 G(S)\longrightarrow Y(S),
\]
is surjective.

The first assertion was proved in \SGARefLink{sga3:xxvi:corollary:5-2}{5.2}.  For the second, let
\(y\in Y(S)\).  Its canonical image in \(X(S)\) is the image of some
\(g\in G(S)\).  The image \(y'\) of \(g\) in \(Y(S)\) therefore has
the same image in \(X(S)\) as \(y\).  There is consequently a unique
\[
 u\in\rad^u(P)(S)
\]
such that \(y'u=y\), and the image of \(gu\) in \(Y(S)\) is indeed
\(y\).

\begin{sgastatement}{\SGAInlineRefTarget{sga3:xxvi:corollary:5-10}Corollary 5.10}
Let \(S\) be a semilocal scheme and \(G\) an \(S\)-reductive group.
\begin{enumerate}[label=\textup{(\roman*)}]
\item Let \(P\) be a parabolic subgroup of \(G\), and \(L\) a Levi
subgroup of \(P\).  The canonical maps \textup{(cf.\ 3.21)}
induce bijections
\[
 H^1(S,L)\xrightarrow{\ \sim\ }H^1(S,P)
       \xrightarrow{\ \sim\ }H^1_{t(P)}(S,G).
\]
\item If
\[
 t,t'\in\Of\bigl(\operatorname{Dyn}(G)\bigr)(S),
\]
then \textup{(cf.\ 3.21)}
\[
 H^1_t(S,G)\cap H^1_{t'}(S,G)=H^1_{t\cap t'}(S,G).
\]
\item If \(P\) and \(Q\) are two parabolic subgroups of \(G\) in
osculating position, the following canonical diagram is cartesian
and all its arrows are injective:
\[
\begin{tikzcd}[column sep=huge,row sep=large]
 H^1(S,P) \arrow[r] & H^1(S,G)\\
 H^1(S,P\cap Q) \arrow[u] \arrow[r] & H^1(S,Q). \arrow[u]
\end{tikzcd}
\]
\end{enumerate}
\end{sgastatement}
\SGARefTarget{sga3:xxvi:cor:5:10}

Let us prove \textup{(i)}.  The map
\[
 H^1(S,L)\longrightarrow H^1(S,P)
\]
is bijective by \SGARefLink{sga3:xxvi:cor:2:3}{2.3}, while
\[
 H^1(S,P)\longrightarrow H^1_{t(P)}(S,G)
\]
is surjective by 3.21.  We must prove that it is injective, or
equivalently that the canonical map
\[
 H^1(S,P)\longrightarrow H^1(S,G)
\]
is injective.  Let \(Q\) be a principal homogeneous bundle under
\(P\), let \(Q_1\) be the associated principal homogeneous bundle
under \(G\), and let \(P'\) and \(G'\) be the corresponding twisted
forms of \(P\) and \(G\).  Clearly \(G'\) is an \(S\)-reductive group
and \(P'\) is a parabolic subgroup of \(G'\).

The elements of \(H^1(S,P)\) having the same image in \(H^1(S,G)\)
as the class of \(Q\) are naturally identified with the kernel of
\[
 H^1(S,P')\longrightarrow H^1(S,G').
\]
By the exact sequence in cohomology, this kernel is in turn
identified with the set of \(G'(S)\)-orbits in
\((G'/P')(S)\).\footnote{For these arguments in nonabelian
cohomology, see Giraud's thesis, \textup{[Gi71]}, \S~III.3.}
But \(G'(S)\) acts transitively on \((G'/P')(S)\), by \SGARefLink{sga3:xxvi:corollary:5-2}{5.2}.

For \textup{(ii)}, let \(Q\) be a principal homogeneous bundle under
\(G\), and let \(G_Q\) be the corresponding twisted form of \(G\).
By definition \textup{(3.21)}, we must prove that \(G_Q\) has a
parabolic subgroup of type \(t\cap t'\) if and only if it has
parabolic subgroups of types \(t\) and \(t'\).  This is precisely the
combination of \SGARefLink{sga3:xxvi:lemma:3:8}{3.8} and \SGARefLink{sga3:xxvi:corollary:5-7}{5.7}\textup{(i)}.  Finally, \textup{(iii)}
follows immediately from \textup{(i)} and \textup{(ii)}.

We now state a consequence of \SGARefLink{sga3:xxvi:cor:5:3}{5.3}.

\begin{sgastatement}{Corollary 5.11}
Let \(S\) be a semilocal scheme, \(G\) an \(S\)-reductive group, and
\(P\) a parabolic subgroup of \(G\).  If \(P\ne G\), there exist at
least three distinct parabolic subgroups of \(G\) having the same
type as \(P\).  In other words,
\[
 P\ne G\qquad\Longrightarrow\qquad
 \bigl(G(S):P(S)\bigr)\geq 3.
\]
\end{sgastatement}
\SGARefTarget{sga3:xxvi:cor:5:11}

Indeed, let \(P'\) be a parabolic subgroup opposite to \(P\), as in
\SGARefLink{sga3:xxvi:corollary:4-3-5}{4.3.5}\textup{(i)}.  Since \(P\ne G\), one has
\(\rad^u(P')\ne e\), for example by \SGARefLink{sga3:xxvi:item:4:3:1}{4.3.2}.  By \SGARefLink{sga3:xxvi:proposition:2-1}{2.1},
\(\rad^u(P')(S)\ne e\); choose
\[
 u\in\rad^u(P')(S),\qquad u\ne e.
\]
Then \(\Int(u)P\ne P\).  By \SGARefLink{sga3:xxvi:cor:5:3}{5.3} there is a parabolic subgroup \(P_1\)
of the same type as \(P\), opposite to both \(P\) and \(\Int(u)P\).
Thus \(P_1\), \(P\), and \(\Int(u)P\) are three distinct parabolic
subgroups of \(G\), all of the same type.
\SGAQED
